\section{对偶变量是约束的影子价格}
\label{sec:09-Convex-Optimization/05-Duality/04}

推理服务的容量规划刚跑完，最优月成本 42 万，在当前的 GPU 配额、显存上限和 \(p99\) 延迟指标下已经压不动了。你把结果发出去，平台组回了一句：再给你八张卡，成本能降多少？降得少就不采购了。

最优解回答不了这个问题。它描述的是配额固定时的最佳做法，配额一改，整个解都得重算。但求解器返回的东西里已经含着答案——那组你可能顺手忽略掉的乘子。

\subsection{把右端项变成自变量，最优值就成了一条曲线}

设想把每条约束的右端松动一点：不等式允许放宽 \(u_i\)，等式的右端挪动 \(v_j\)。最优值随之变成扰动量的函数

\[
p^\star(u,v)=\inf\set{f_0(x)\given f_i(x)\le u_i,\ h_j(x)=v_j}.
\]

零扰动处 \(p^\star(0,0)\) 就是原来那个 42 万。

扰动后的 Lagrangian 只比原来多两个不含 \(x\) 的平移项，

\[
L(x,\lambda,\nu;u,v)=f_0(x)+\lambda^{\trans}\bigl(f(x)-u\bigr)+\nu^{\trans}\bigl(h(x)-v\bigr),
\]

对 \(x\) 取下确界时这两项直接提到外面：

\[
\inf_x L(x,\lambda,\nu;u,v)=g(\lambda,\nu)-\lambda^{\trans}u-\nu^{\trans}v,
\]

其中 \(g\) 是原来那个未扰动问题的对偶函数。注意对偶可行性只要求 \(\lambda\ge0\)，与 \(u,v\) 无关，所以原问题的最优乘子 \((\lambda^\star,\nu^\star)\) 对扰动后的问题仍然合法，弱对偶立刻给出

\[
p^\star(u,v)\ge g(\lambda^\star,\nu^\star)-{\lambda^\star}^{\trans}u-{\nu^\star}^{\trans}v.
\]

再假设未扰动问题满足强对偶，右端第一项就等于 \(p^\star(0,0)\)：

\[
p^\star(u,v)\ge p^\star(0,0)-{\lambda^\star}^{\trans}u-{\nu^\star}^{\trans}v.
\]

这是一个对所有 \(u,v\) 成立的仿射下界，而且在 \(u=v=0\) 处取到等号。一条曲线被一个仿射函数从下方托住并在某点相切，那个仿射函数的斜率就是曲线在该点的斜率（可微时）或次梯度（一般情形）。于是

\[
\frac{\partial p^\star}{\partial u_i}(0)=-\lambda_i^\star,
\qquad
\frac{\partial p^\star}{\partial v_j}(0)=-\nu_j^\star.
\]

乘子的身份到这里就定下来了：它是最优值对约束右端项的敏感度，差一个负号。负号来自最小化——放宽上界让可行集变大，成本只会降不会升。改成最大化利润，符号翻过来。

\begin{center}
\begin{tikzpicture}[x=1cm,y=1cm,font=\small,>=stealth]
  \draw[->,black!60] (0,0) -- (6.8,0) node[right,font=\footnotesize] {约束放松量 \(u\)};
  \draw[->,black!60] (0.2,0) -- (0.2,3.3) node[above,font=\footnotesize] {\(p^\star(u)\)};
  \draw[dashed,black!40] (2.0,0) -- (2.0,1.60);
  \draw[dashed,black!40] (0.2,1.60) -- (2.0,1.60);
  \node[font=\footnotesize,black!60,below] at (2.0,-0.05) {\(u=0\)};
  \draw[thick,red!60!black] (0.6,2.51) -- (4.2,0.17);
  \draw[very thick,blue!55!black] plot[smooth] coordinates
    {(0.5,2.80) (1.0,2.35) (1.5,1.95) (2.0,1.60) (2.5,1.30) (3.0,1.05)
     (3.5,0.85) (4.0,0.70) (4.5,0.62) (5.0,0.56) (5.5,0.52) (6.0,0.50)};
  \draw[dashed,black!55] (2.0,1.60) -- (3.0,1.60) -- (3.0,0.95);
  \node[font=\footnotesize,black!70,above] at (2.5,1.62) {\(1\)};
  \node[font=\footnotesize,black!70,right] at (3.05,1.27) {\(\lambda^\star\)};
  \node[font=\footnotesize,red!60!black,right] at (4.1,0.25) {切线斜率 \(=-\lambda^\star\)};
  \draw[fill=black!70] (4.0,0.70) circle (0.06);
  \node[font=\footnotesize,black!65,align=left,right] at (4.15,0.95)
    {活跃集在此换挡，\\价格随之改变};
  \node[font=\footnotesize,blue!55!black,left] at (0.95,3.0) {};
\end{tikzpicture}

\emph{最优值随约束右端项变化的曲线；它在 \(u=0\) 处的斜率就是 \(-\lambda^\star\)，虚线三角形读作``多给一个单位，成本降 \(\lambda^\star\)''。}
\end{center}

图里那条切线在曲线下方，这是凸性的直接后果，也是一句实用的提醒：按乘子估算出来的收益偏乐观。放松一个单位省 \(\lambda^\star\)，放松十个单位省的不到 \(10\lambda^\star\)，因为曲线一路变平。

\subsection{量纲让价格变得可以核对}

设目标的单位是元，某条约束写成 \(a^{\trans}x\le b\)，\(b\) 的单位是 GPU 卡时。Lagrangian 里出现的是 \(\lambda(a^{\trans}x-b)\)，这一项必须与目标同量纲，所以 \(\lambda\) 的单位是元每卡时。它可以直接与市场报价比大小：租一张卡的实际单价低于 \(\lambda\) 就该租，高于 \(\lambda\) 就不该。开头那个八张卡的问题，答案是 \(8\lambda^\star\) 与采购成本的比较，前提是八张卡还落在切线适用的范围内。

量纲还能抓错。把约束按百分比重写，右端缩小 100 倍，乘子就自动放大 100 倍；乘子的数值随约束的写法变，但``乘子乘以右端变化量''这个乘积始终是元，不受缩放影响。跨问题比较乘子数值之前，先确认两边的约束是按同一种单位写的。

\subsection{正则化系数就是范数预算的影子价格}

这层含义在建模里出现得比在工厂里还频繁。带范数约束的经验风险最小化写成

\[
\begin{aligned}
\text{minimize}\quad&L(w)\\
\text{subject to}\quad&\norm{w}\le C,
\end{aligned}
\]

它的 Lagrangian 是 \(L(w)+\lambda(\norm{w}-C)\)。固定住最优乘子 \(\lambda^\star\)，常数项 \(-\lambda^\star C\) 不影响极小点，于是求解上面这个约束问题与求解

\[
\text{minimize}\quad L(w)+\lambda^\star\norm{w}
\]

得到同一个 \(w^\star\)。日常写的罚项形式与预算形式，是同一个问题的两副写法，\(\lambda\) 与 \(C\) 通过对偶一一对应。

对应关系里几件事值得记住。第一，\(\lambda^\star=-\mathrm{d}p^\star/\mathrm{d}C\)：调大 \(\lambda\) 等价于调小预算 \(C\)，两个旋钮方向相反。第二，互补松弛在这里给出一个干净的判据——若无约束的极小点本身就满足 \(\norm{w}<C\)，约束是松的，\(\lambda^\star=0\)，正则项一点不该加；正则化起作用的前提是预算确实卡住了。第三，\(C\mapsto\lambda^\star(C)\) 这个映射依赖数据，事先算不出来，所以实践中调的是 \(\lambda\) 而不是 \(C\)，尽管后者在业务上更好解释。给非专业听众汇报时，把选定的 \(\lambda\) 反查成对应的 \(\norm{w^\star}\)，比报一个裸的超参数值有说服力得多。\(L_1\) 与 \(L_2\) 两种罚项各自的几何后果见\hyperref[sec:09-Convex-Optimization/07-Applications/02]{``\(L_1\) 正则为何产生稀疏''}与\hyperref[sec:09-Convex-Optimization/07-Applications/01]{``光滑经验风险与 \(L_2\) 正则化''}。

\subsection{先看哪些约束有价，再看放松哪个划算}

一组乘子解出来之后，最有操作性的读法是把约束按 \(\lambda_i^\star\) 排个序。乘子为零的约束不必碰——互补松弛说它们还有余量，多给资源改变不了任何东西。乘子为正的才是瓶颈，而瓶颈之间的排序要用 \(\lambda_i^\star\) 除以放松一个单位的真实代价：一条约束的影子价格是每卡时 3 元而卡时的采购价是 2 元，另一条是每 GB 显存 0.4 元而扩显存的边际成本是 0.6 元，该动的是前者。

这套读法有边界，报告之前要说清楚。放松一条约束，最优解会沿某个方向滑动，滑到一定程度会撞上另一条原本松着的约束，活跃集一变，斜率就换挡——图里那个黑点标的正是这件事，价格只在换挡之前有效。乘子还可能不唯一：几条活跃约束的梯度线性相关时，满足 KKT 的乘子构成一个集合，对应的是价值函数在该点的次微分而非一个数，此时单独一个乘子没有定价含义，只有它们与约束变动的组合才有。这种退化在\hyperref[sec:09-Convex-Optimization/05-Duality/03]{``KKT 把最优性写成一组方程''}里已经出现过，那里表现为约束取等而乘子为零，这里表现为价格不唯一，来源是同一个。

还有一条前提容易被跳过：上面全部推导用了强对偶。间隙不为零时，那个仿射下界依然成立，但它不再在 \(u=0\) 处相切，斜率与真实的敏感度对不上。整数规划松弛给出的乘子常被当作影子价格用，用之前要清楚它描述的是松弛问题的敏感度，与原问题之间隔着一条整数间隙。

\subsection{对偶给的是解释，不只是解}

把本单元的四篇串起来看，对偶提供的是同一件东西的四种用法。乘子构造出全局下界，下界与任何可行解一夹就成为最优性证书；间隙是否闭合取决于凸性与约束资格，闭不上时下界仍然可用，只是不会自己收敛到零；间隙闭合的那一刻，最优性坍缩成 KKT 那组方程，其中互补松弛把约束分成有价与无价两类；而乘子的数值本身，是最优值对约束右端的敏感度。

最后这一层把对偶从求解技术变成了决策语言。原始解回答的是该怎么做，对偶解回答的是为什么只能做到这里、哪条限制在卡着、松开它值不值。一份只有原始解的报告没法被审查，因为读者无从判断哪个约束是关键的；带上乘子，这份报告就同时给出了行动方案和它的定价依据。

下一单元转向算法本身。届时会看到内点法在中心路径上同时推进原始变量与乘子，增广 Lagrange 与交替方向法把乘子当作协调子问题的价格信号在迭代中更新——这些方法的收敛判据，正是这里的间隙与 KKT 残差，见\hyperref[ch:076]{``优化算法：怎样把局部模型变成全局进展''}。
